\clearpage\pagestyle{otMain}\setlength{\otBlocEspacement}{2.5mm}\fontsize{11}{15}\selectfont\corrigesChapitre{2}{Repérer, comparer et arrondir}
\begin{blocLivre}{cor-2-1}\begin{corrige}{1}{Lire des abscisses}\demarche{Trouver la valeur d’un petit intervalle.}
Chaque petit intervalle vaut $0{,}1$. A : $\res{0{,}6}$ ; B : $\res{1{,}3}$ ; C : $\res{2{,}8}$.
\label{fig-cor-2-1}\figureGeo[corrige]{rep-1}
\end{corrige}\end{blocLivre}
\begin{blocLivre}{cor-2-2}\begin{corrige}{2}{Placer des points}\demarche{Repérer le dixième, puis compter les centièmes.}
D est 5 centièmes après $5{,}2$ ; E est 8 centièmes après $5{,}3$ ; F est exactement sur la graduation $5{,}3$. Les positions figurent en vert.
\label{fig-cor-2-2}\figureGeo[corrige]{rep-2}
\end{corrige}\end{blocLivre}
\begin{blocLivre}{cor-2-3}\begin{corrige}{3}{Comparer}\demarche{Comparer partie entière, puis dixièmes, puis centièmes.}
$\res{12{,}7>12{,}69}$ ; $\res{0{,}8>0{,}75}$ ; $\res{3{,}05<3{,}5}$ ; $\res{7>6{,}999}$.
\end{corrige}\end{blocLivre}
\begin{blocLivre}{cor-2-4}\begin{corrige}{4}{Ranger}\demarche{Écrire tous les nombres avec le même nombre de décimales.}
Avec trois décimales : $4{,}200$ ; $4{,}020$ ; $4{,}250$ ; $4{,}205$ ; $4{,}300$. Ordre : $\res{4{,}02<4{,}2<4{,}205<4{,}25<4{,}3}$.
\end{corrige}\end{blocLivre}
\begin{blocLivre}{cor-2-5}\begin{corrige}{5}{Arrondir}\demarche{Regarder le chiffre qui suit le rang voulu.}
$7{,}364$ : $\res7$ ; $\res{7{,}4}$ ; $\res{7{,}36}$. $12{,}58$ : $\res{13}$ ; $\res{12{,}6}$ ; $\res{12{,}58}$, qui a déjà deux décimales.
\end{corrige}\end{blocLivre}
\begin{blocLivre}{cor-2-6}\begin{corrige}{6}{Encadrer}\demarche{Garder les chiffres jusqu’au rang voulu, puis ajouter une unité de ce rang.}
$\res{15<15{,}738<16}$ ; $\res{15{,}7<15{,}738<15{,}8}$ ; $\res{15{,}73<15{,}738<15{,}74}$.
\end{corrige}\end{blocLivre}
\begin{blocLivre}{cor-2-7}\begin{corrige}{7}{Intercaler}\demarche{Écrire les bornes avec une décimale de plus.}
Par exemple $\res{2{,}35}$ ; $\res{0{,}595}$ ; $\res{8{,}005}$. D’autres réponses conviennent : il suffit d’ajouter un rang.
\end{corrige}\end{blocLivre}
\begin{blocLivre}{cor-2-8}\begin{corrige}{8}{Zoom sur la droite}\demarche{Un demi-centième vaut 5 millièmes.}
G est 7 centièmes après $1{,}4$ : $\res{1{,}47}$. H est au milieu entre $1{,}41$ et $1{,}42$ : $\res{1{,}415}$, soit 5 millièmes de plus que $1{,}41$.
\label{fig-cor-2-8}\figureGeo[corrige]{rep-8}
\end{corrige}\end{blocLivre}
\begin{blocLivre}{cor-2-9}\begin{corrige}{9}{Valeurs approchées}\demarche{Arrondir comme pour un décimal.}
$\frac13$ : $\res{0{,}3}$ et $\res{0{,}33}$. $\frac23$ : $\res{0{,}7}$ et $\res{0{,}67}$, car le chiffre suivant est 6.
\end{corrige}\end{blocLivre}
\begin{blocLivre}{cor-2-10}\begin{corrige}{10}{Tous les candidats}\demarche{Chercher les bornes, chiffre 5 compris.}
Ils vont de $3{,}45$ à $3{,}54$ : $3{,}45$ ; $3{,}46$ ; … ; $3{,}54$. Il y en a $\res{10}$. $3{,}44$ s’arrondit à $3{,}4$ et $3{,}55$ à $3{,}6$.
\end{corrige}\end{blocLivre}
\begin{blocLivre}{cor-2-11}\begin{corrige}{11}{Le saut en longueur}\demarche{Comparer avec le même nombre de décimales.}
Avec deux décimales : $4{,}90$ ; $4{,}89$ ; $4{,}85$ ; $4{,}80$ ; $4{,}08$. Podium : $\res{\text{Malik, Tom, Léna}}$ ; puis Inès et Sofia.
\end{corrige}\end{blocLivre}
\begin{blocLivre}{cor-2-12}\begin{corrige}{12}{Les prix arrondis}\demarche{Arrondir pour estimer, puis calculer exactement.}
Arrondis : 2 € ; 14 € ; 1 €. Estimation : $\res{17}$ €. Prix exact : $2{,}49+13{,}95+0{,}99=\res{17{,}43}$ €, proche de l’estimation.
\end{corrige}\end{blocLivre}
\begin{blocLivre}{cor-2-13}\begin{corrige}{13}{Le chronomètre}\demarche{Comparer les valeurs exactes, puis les arrondis.}
Ordre : Bilal $12{,}400$ ; Dan $12{,}405$ ; Ana $12{,}480$ ; Chloé $12{,}520$ : $\res{\text{Bilal gagne}}$. Au dixième : Ana $12{,}5$ ; Bilal $12{,}4$ ; Chloé $12{,}5$ ; Dan $12{,}4$. $\res{\text{Non}}$ : Bilal et Dan seraient à égalité, tout comme Ana et Chloé.
\end{corrige}\end{blocLivre}
\begin{blocLivre}{cor-2-14}\begin{corrige}{14}{Le nombre mystère}\demarche{Utiliser l’arrondi le plus précis d’abord.}
Arrondi au dixième $7{,}6$ : le nombre est entre $7{,}55$ et $7{,}64$. Avec 9 centièmes, seul $\res{7{,}59}$ convient ; son arrondi à l’unité est bien 8.
\end{corrige}\end{blocLivre}
\begin{blocLivre}{cor-2-15}\begin{corrige}{15}{Vrai ou faux}\demarche{Comparer rang par rang.}
C’est $\res{\text{faux}}$ : les parties entières sont égales, puis 1 dixième est moins que 9 dixièmes. Donc $\res{3{,}15<3{,}9}$, comme le montre l’écriture $3{,}15<3{,}90$.
\end{corrige}\end{blocLivre}
